\section{Vector V2: Input Data Manipulation}
\label{sec:v2}

\subsection{Attack Description}

The \bisect{} algorithm ingests three data sources: the Census P.L.\ 94-171
redistricting file (population counts), TIGER/Line shapefiles (geographic
boundaries), and optionally American Community Survey 5-year tables
(demographic composition).
A hostile actor with control over the data pipeline could modify population
counts, tract boundaries, or demographic tables before the algorithm
processes them.

The three attack surfaces are:
\begin{enumerate}
\item \textbf{Population inflation/deflation}: Systematically inflate tract
  populations in areas favorable to one party and deflate in areas favorable
  to the other.
  This biases the equal-population constraint toward partisan outcomes.
\item \textbf{Shapefile manipulation}: Modify tract boundary coordinates to
  change adjacency relationships, effectively rewiring the graph without
  changing population counts.
\item \textbf{ACS table manipulation}: Alter minority population percentages
  used in VRA boost weighting to weaken or strengthen minority community
  protections in targeted areas.
\end{enumerate}

\subsection{Maximum Achievable Shift}

\paragraph{Population manipulation.}
A $\pm 1\%$ population shift per tract, applied strategically to all tracts
in competitive states, produces at most $\Delta D \leq 0.3$ seats nationally.
This bound comes from the observation that the algorithm's equal-population
constraint ensures that any distortion propagates across multiple districts
within a state, partially canceling the partisan effect.
Furthermore, a $\pm 1\%$ distortion is within the margin of B.7's seed
noise floor and would be essentially undetectable without the SHA-256 chain.

Larger distortions (e.g., $\pm 5\%$) would produce larger effects
but are detectable at high probability through the verification chain
(Section~\ref{sec:v2-detection}).

\paragraph{Shapefile manipulation.}
Modifying tract boundaries to alter adjacency relationships is the most
consequential V2 attack.
Merging two adjacent tracts of different partisan compositions into a
single node, then splitting them differently, can effectively pre-gerrymander
the graph before the algorithm runs.
We estimate this attack can produce at most $\Delta D \leq 0.4$ seats
nationally before the adjacency changes become geometrically implausible
(creating non-contiguous tracts or implausible boundary shapes).

\paragraph{ACS manipulation.}
VRA boost weighting affects compactness slightly but has negligible partisan
effect ($\Delta D \leq 0.1$ nationally from B.17's $w_{\rm vra}$ sweep).
ACS manipulation thus does not constitute a meaningful partisan gaming vector;
it is primarily a VRA compliance concern.

\subsection{Detection Mechanism}
\label{sec:v2-detection}

The \bisect{} fetch and verification chain provides three levels of defense:
\begin{enumerate}
\item \textbf{Download hash}: \texttt{bisect fetch} computes a SHA-256 hash
  of each downloaded Census P.L.\ 94-171 file and TIGER/Line shapefile
  immediately upon download.
  These hashes are stored in the manifest and compared against the Census
  Bureau's published checksums (available at
  \url{https://www2.census.gov/}).
\item \textbf{Adjacency hash}: \texttt{bisect} computes a hash of the
  adjacency graph constructed from the TIGER/Line shapefiles.
  Any modification to shapefile boundaries that changes adjacency
  relationships will change the adjacency hash.
\item \textbf{Plan verification}: \texttt{bisect label-verify} replays the
  entire computation from the stored hashes and checks that the output
  plan's hash matches the manifest.
  Any inconsistency in the chain---from raw Census file through adjacency
  construction to final plan---is flagged.
\end{enumerate}

\paragraph{Security limitation.}
\bisect{} uses TLS certificate validation via the system trust store for
Census Bureau downloads but does \emph{not} implement certificate pinning.
This is a disclosed limitation: a sophisticated man-in-the-middle attacker
with the ability to spoof TLS certificates could intercept downloads.
In practice, the Census Bureau's P.L.\ 94-171 files are also available
on physical media and through state redistricting data services, providing
alternative verification paths.
For maximum security, practitioners should verify download hashes against
multiple independent sources.

Detection probability under normal \dia{} deployment: $\delta_2 \approx 1.00$
for population and ACS manipulation; $\delta_2 \approx 0.98$ for shapefile
manipulation (bounded by the fraction of adjacency changes that also produce
geometrically implausible boundaries).
